{\fontsize{9.5}{12}\selectfont\renewcommand{\arraystretch}{1.13}
\begin{longtable}{@{}P{43mm}P{37mm}P{72mm}@{}}
\toprule 文件及性质 & 定位 & 对本项目的适用要点\\\midrule\endfirsthead
\toprule 文件及性质 & 定位 & 对本项目的适用要点\\\midrule\endhead
\bottomrule\endlastfoot
生成式人工智能服务管理暂行办法（部门规章）\src{REG-001} & 第2、4(5)、17条 & 先判断是否向境内公众提供服务；透明、准确可靠及备案／评估义务须结合适用范围和条件。内部应用排除条款不免除其他法律责任。\\
个人信息保护法（法律）\src{REG-002} & 第13、21、24、27、54--56条 & 多种合法基础；委托约定与监督；自动化决策的透明、公平和特定说明／拒绝权；公开信息合理范围；定期审计及特定处理事前影响评估。评估报告和处理记录至少保存3年。\\
证券基金经营机构信息技术管理办法（证监会规章，2021修订）\src{REG-003} & 第3、43--45、63条 & 按实际IT服务角色判断；保留自身责任和重要系统控制，外部独立运维限制有法定例外；审查报送、替换预案及协议。第63条定义重要信息系统。\\
网络数据安全管理条例（行政法规，2025施行）\src{REG-004} & 第9、12、19、29、31、33、37条 & 安全措施、委托／提供个人信息与重要数据的合同监督及处理记录；训练数据规则结合生成服务角色理解；重要数据按目录、告知及识别义务判断。\\
证券期货业网络和信息安全管理办法（令218）\src{REG-005} & 第2、4、29--34、71、75条 & 证券基金机构网络信息安全及投资者信息保护责任；第75条废止旧信息安全保障办法。\\
促进和规范数据跨境流动规定（部门规章）\src{REG-006} & 第5(4)、7--8、10--11、13条 & 出境机制按关基／重要数据和累计人数等分层；普通与敏感个人信息门槛不同；豁免机制仍保留合法处理等义务。需公司级事实，不能按原型样本量代替累计统计。\\
网络安全法（2025修正，2026施行）\src{REG-013} & 第23条 & 等保相关措施、网络运行和安全事件监测记录；不少于6个月的期限针对相关网络日志，不能直接等同全部提示词原文保留期。\\
数据安全法（法律）\src{REG-014} & 第21、27、29--32条 & 数据分类分级、全流程管理、风险监测、重要数据风险评估和合法收集。保护对象不限于个人信息。\\
个人信息保护合规审计管理办法（2025施行）\src{REG-015} & 第3--4条 & 定期开展合规审计；处理超过1,000万人个人信息的，至少每两年一次。低于门槛不等于免审计。\\
人工智能生成合成内容标识办法（2025施行）\src{REG-016} & 第2--6、9条 & 按生成服务、传播平台、发布者及具体内容类型判断显式／隐式标识与用户声明义务；对外发布需重判角色。配套强制国标技术实现需在架构确定后核对。\\
基金经营机构大模型技术应用规范（团体标准）\src{REG-009}\src{REG-010} & 9.1.2、9.2、11.2--11.3；附录A／B & 自有场景评测、部署、数据权限与解释参考；11.3(e)用“宜”，附录为资料性。适用依据还须结合标准化法第18条与实际采用约定。\src{REG-012}\\
SR 26-2及配套银行MRM指引（美国监督指引）\src{REG-007}\src{REG-008} & 发布函；指引I--II节、脚注3 & 2026-04-17替代旧指引；非可强制执行标准，面向所述银行机构。生成式／Agentic AI不在定义范围，机构仍自行确定风险治理。非MSIM内部政策。\\
\end{longtable}}
